\abstract{Access-to-information law presumes that public administrations hold an account of how
their procedures actually run. Empirical work on open-records requests finds that this account
is frequently absent: the binding obstacle is not exemption but that the record was never
created. We derive such a record from an administration's own event log. From 5{,}243 concluded
administrative cases at a state-level secretariat of a Brazilian state administration, we build
an event knowledge graph, group cases into emergent process types from opening-document text
together with routing and action history, and learn a typical path for each type using
sequence models over two trace alphabets: recorded actions, and the itinerary of organisational
units a case visits. The catalogue is reproducible --- every clustering refittable from a feature
matrix reproduces exactly from its stamped features, every reported table is regenerated from
artefacts, and a register of the load-bearing figures in the running text is checked against
those artefacts in both directions. Evaluating it required controls that this
literature does not routinely report. A shuffled partition, preserving cluster count, size
distribution and evaluated population, reduces next-step accuracy from 0.726 to 0.494,
establishing that the grouping carries signal. What per-cluster modelling then contributes
depends on how it is measured, and the conventional measurement overstates it roughly sixfold.
Against the usual within-cluster baseline the apparent gain on action prediction is 0.083;
against a single global model with the model family held fixed it is 0.004; and once both trace
alphabets are scored on one population it is negative, between $-0.009$ and $-0.005$ depending
on whether clusters are weighted by membership or by prediction steps. Routing gains 0.029 to
0.036 on the same population and under the same choice. Two further results constrain the
method: sixteen clusterings across four algorithm families fall within 0.0079 once scored on a
common set of cases, on an interval too wide to call them equivalent; and the internal validity
indices that would ordinarily choose between them track cluster granularity rather than
downstream utility, correlating at 0.87 with median cluster size while their association with
accuracy cannot be resolved at this sample size. Anonymisation applied before clustering, as the
governing instrument requires, costs no measurable next-activity utility against a pre-registered
margin.}

\keywords{process mining, administrative transparency, trace clustering, access to information,
data protection, public administration}
